\newpage
\section{Hiện thực plugin trong Paint.NET}

\subsection{Nền tảng công nghệ}

Plugin được xây dựng bằng C\# trên nền .NET, hướng tới môi trường Windows với \texttt{TargetFramework} là \texttt{net9.0-windows}. Ở phía Paint.NET, effect được hiện thực thông qua \texttt{BitmapEffect} API thay vì mô hình \texttt{PropertyBasedBitmapEffect} cũ. Cách tiếp cận này cho phép đồ án xây dựng một hộp thoại cấu hình riêng bằng WinForms, từ đó kiểm soát trực tiếp trải nghiệm preview và thời điểm áp dụng hiệu ứng lên canvas.

Ở mức thấp hơn, plugin thao tác với dữ liệu ảnh dưới dạng \texttt{IBitmapLock} và mảng \texttt{ColorBgra32}. Cách tổ chức này giúp quá trình tính toán DCP bám sát bộ nhớ ảnh gốc, hạn chế lớp trừu tượng không cần thiết trong một bài toán phụ thuộc mạnh vào truy cập theo pixel.

\subsection{Cấu trúc source code}

Mã nguồn hiện tại được chia thành các khối chức năng chính:
\begin{itemize}
\item \texttt{.csproj}: khai báo target framework, tham chiếu thư viện Paint.NET và cấu hình build trên Windows;
\item \texttt{DcpDehazeEffect.cs}: điểm vào của plugin, chịu trách nhiệm nhận token cấu hình, tạo form, quản lý cache full render và xuất dữ liệu ra \texttt{IBitmapEffectOutput};
\item \texttt{DehazeProcessor.cs}: lõi xử lý ảnh, chứa các bước như xây dark channel, ước lượng atmospheric light, guided filter, tạo \texttt{omegaMap} và khôi phục ảnh cuối;
\item \texttt{DehazeConfigForm.cs} và \texttt{DehazeConfigToken.cs}: hiện thực giao diện cấu hình và đối tượng truyền trạng thái giữa form với effect;
\item \texttt{DehazeCLI.cs}: tiện ích chạy thuật toán ngoài môi trường Paint.NET để kiểm thử độc lập.
\end{itemize}

Việc tách lõi xử lý khỏi lớp effect giúp cùng một pipeline có thể tái sử dụng cho preview, render toàn ảnh và cả công cụ dòng lệnh. Đây là điểm quan trọng để plugin vừa dễ kiểm thử, vừa thuận tiện cho việc mở rộng sau này.

\subsection{Giao diện cấu hình}

Giao diện của plugin được thiết kế theo hướng tối giản nhưng đủ cho thao tác thực tế. Người dùng có:
\begin{itemize}
\item thanh trượt \texttt{Strength} trong miền \(0\) đến \(100\);
\item ô chọn \texttt{Adaptive haze removal};
\item vùng preview dạng \texttt{PictureBox};
\item hai nút \texttt{OK} và \texttt{Cancel}.
\end{itemize}

Thanh \texttt{Strength} điều khiển mức can thiệp của thuật toán, trong khi tùy chọn \texttt{Adaptive} quyết định có dùng \texttt{omegaMap} theo pixel hay không. Khi hộp thoại đang mở, plugin chỉ cập nhật preview; ảnh trên canvas chính chỉ được render lại khi người dùng bấm \texttt{OK}. Nhờ đó, thao tác tinh chỉnh tham số không làm gián đoạn trải nghiệm sử dụng.

Về mặt hiện thực, khi \texttt{Strength} thay đổi, hệ thống không tính lại toàn bộ pipeline trên ảnh đầy đủ mà chỉ render lại vùng preview từ cache đã được dựng sẵn. Điều này giúp hộp thoại phản hồi nhanh hơn trong quá trình kéo thanh trượt.

\begin{figure}[H]
  \centering
  \includegraphics[width=0.75\textwidth]{plugin-ui.jpg}
  \caption{Giao diện sử dụng effect DHUD DCP Dehaze với thanh Strength và tùy chọn Adaptive haze removal}
\end{figure}

\subsection{Preview và cache}

Một điểm then chốt trong hiện thực là cơ chế preview thu nhỏ. Khi form cấu hình được mở, plugin đọc toàn bộ layer nguồn, resize ảnh sao cho cạnh dài nhất không vượt quá \(320\) pixel, rồi xây trước một cache dehaze trên ảnh thu nhỏ. Sau đó, mỗi lần \texttt{Strength} thay đổi, preview chỉ cần render lại từ cache thay vì tính lại toàn bộ dark channel, atmospheric light và transmission.

Ở mức render toàn ảnh, cache trong \texttt{DcpDehazeEffect} cũng được tổ chức riêng và chỉ được tái sử dụng khi đồng thời thỏa các điều kiện sau:
\begin{itemize}
\item cùng document;
\item cùng source layer;
\item cùng kích thước ảnh;
\item cùng chế độ adaptive.
\end{itemize}

Cách tổ chức này tách rõ phần tiền xử lý nặng khỏi phần điều chỉnh tham số nhẹ. Những bước như tạo dark channel, ước lượng \(A\), dựng transmission và \texttt{omegaMap} chỉ cần thực hiện lại khi đầu vào thực sự thay đổi.

\begin{lstlisting}[style=csharpblock,caption={Trích đoạn kiểm tra điều kiện tái sử dụng cache}]
if (this.cache is { } existing &&
    existing.Adaptive == this.adaptive &&
    ReferenceEquals(existing.Document, currentDocument) &&
    existing.SourceLayerIndex == currentSourceLayerIndex &&
    existing.Data.Width == size.Width &&
    existing.Data.Height == size.Height)
{
    return existing;
}
\end{lstlisting}

\subsection{Quy trình build và tích hợp}

Trên Windows, dự án được build như một thư viện plugin cho Paint.NET. Sau khi chạy \texttt{dotnet build} hoặc \texttt{dotnet build -c Release}, file đầu ra chính là \texttt{DHUD.PaintNet.BitmapCompatibilityEffect.dll} trong thư mục \texttt{bin}. Nếu thiết lập biến \texttt{InstallPlugin=true}, target \texttt{InstallToPaintDotNet} trong file \texttt{.csproj} còn có thể tự chép thư viện vào thư mục \texttt{Effects} của Paint.NET.

Trong trường hợp build thủ công, chỉ cần chép file \texttt{.dll} vào thư mục \texttt{Effects} rồi khởi động lại ứng dụng. Sau khi được nạp, effect xuất hiện dưới tên \texttt{DHUD DCP Dehaze} trong submenu \texttt{DHUD Restoration}. Điều này cho thấy sản phẩm cuối không chỉ dừng ở mức mã nguồn thuật toán, mà đã được đóng gói thành một plugin có thể sử dụng trực tiếp trong quy trình chỉnh sửa ảnh.

\begin{figure}[H]
  \centering
  \includegraphics[width=0.88\textwidth]{plugin_place.jpg}
  \caption{Vị trí của plugin trong Paint.NET: Effects > DHUD Restoration > DHUD DCP Dehaze}
\end{figure}
